\documentclass[10pt]{article}
\usepackage[T1]{fontenc}
\usepackage[utf8]{inputenc}
\usepackage[french]{babel}
\usepackage{amsmath,amssymb,mathrsfs,mathpazo}
\usepackage{geometry}
\geometry{a4paper,margin=2cm}
\usepackage{graphicx,array,tabularx,tabularray,colortbl,multirow,multicol,booktabs,enumitem,pifont,ulem,fancybox,diagbox,cancel}
\usepackage{pstricks,pstricks-add,pst-plot,pst-tree,pst-node,pst-text,pst-func,pst-math,pst-bspline,pst-3dplot,pst-eucl,pst-solides3d}
\let\pscalculate\undefined
\usepackage{pst-bezier}
\usepackage{tikz,pgfplots,tkz-tab}
\usetikzlibrary{arrows,arrows.meta,patterns,calc,positioning,trees,decorations.markings,shapes.geometric,shapes.misc,shapes.arrows,intersections,decorations.pathmorphing,decorations.pathreplacing,backgrounds,shadows,babel}
\pgfplotsset{compat=1.18}
\usepgfplotslibrary{fillbetween}
\usepackage[np]{numprint}
\usepackage{listings,makecell,pas-tableur,twoopt}
\let\barre\undefined
\usepackage[pstricks]{tablvar}
\definecolor{darkgreen}{RGB}{0,100,0}
\providecommand{\red}{\color{red}}
\providecommand{\blue}{\color{blue}}
\providecommand{\green}{\color{green}}
\providecommand{\black}{\color{black}}
\providecommand{\white}{\color{white}}
\providecommand{\gray}{\color{gray}}
\providecommand{\cyan}{\color{cyan}}
\providecommand{\magenta}{\color{magenta}}
\providecommand{\yellow}{\color{yellow}}
\providecommand{\violet}{\color{violet}}
\providecommand{\orange}{\color{orange}}
\providecommand{\euro}{\texteuro}
\providecommand{\EUR}{\texteuro}
\providecommand{\ieme}{\textsuperscript{e}}
\providecommand{\degres}{\ensuremath{^\circ}}
\providecommand{\degree}{\ensuremath{^\circ}}
\providecommand{\vv}[1]{\overrightarrow{#1}}
\providecommand{\Coord}[1]{#1}
\providecommand{\decosix}{}
\providecommand{\starredbullet}{\ensuremath{\bullet}}
\pagestyle{empty}
\setlength{\parindent}{0pt}
\usepackage{adjustbox}\DeclareUnicodeCharacter{2236}{:}\DeclareUnicodeCharacter{2190}{\ensuremath{\leftarrow}}
\providecommand{\R}{\mathbb{R}}
\providecommand{\N}{\mathbb{N}}
\providecommand{\Z}{\mathbb{Z}}
\providecommand{\Q}{\mathbb{Q}}
\providecommand{\D}{\mathbb{D}}
\providecommand{\C}{\mathbb{C}}
\providecommand{\e}{{\mathrm{e}}}
\providecommand{\vect}[1]{\overrightarrow{#1}}
\providecommand{\vectt}[1]{\overrightarrow{\mathrm{#1}}}
\providecommand{\vvt}[1]{\overrightarrow{\mathrm{#1}}}
\providecommand{\np}[2][]{{#2~#1}}
\providecommand{\vv}[1]{{\overrightarrow{#1}}}
\renewcommand{\d}{{\mathrm{d}}}
\providecommand{\degres}{{^\circ}}
\providecommand{\degre}{{^\circ}}
\providecommand{\blue}{\color{blue}}
\providecommand{\red}{\color{red}}
\providecommand{\enskip}{\quad}
\providecommand{\llbracket}{\lbrack\!\lbrack}
\providecommand{\rrbracket}{\rbrack\!\rbrack}
\providecommand{\euro}{{\text{€}}}
\providecommand{\no}{{\text{n°}}}
\providecommand{\interval}[2]{{[#1;#2]}}
\providecommand{\intervFO}[2]{{[#1;#2[}}
\providecommand{\dotfill}{\ldots}

\providecommand{\ds}{\displaystyle}
\providecommand{\textding}[1]{#1}
\providecommand{\Oij}{\(\mathrm{O};\vec{\imath},\vec{\jmath}\)}
\providecommand{\Oijk}{\(\mathrm{O};\vec{\imath},\vec{\jmath},\vec{k}\)}
\providecommand{\Ouv}{\(\mathrm{O};\vec{u},\vec{v}\)}
\providecommand{\intervFF}[2]{\ensuremath{[#1;#2]}}
\renewcommand{\euro}{\ifmmode\text{\texteuro}\else\texteuro\fi}
\DeclareUnicodeCharacter{00B0}{\ifmmode^{\circ}\else\textdegree\fi}
\DeclareUnicodeCharacter{20AC}{\ifmmode\text{\texteuro}\else\texteuro\fi}
% Énoncés : Éducation nationale. Transcription et corrigés : APMEP, Denis Vergès et contributeurs. Découpage et classement : CoopMaths, Sébastien Lozano, Éric Elter et contributeurs. Adaptation de la présentation : Maths Explorer.
% CC BY-SA : https://coopmaths.fr/www/about/
% https://coopmaths.fr/alea/static/bac/2022/tex/bac_2022_05_sujet2_ameriquenord_4_cor.tex
\begin{document}
\noindent\textit{APMEP / CoopMaths. CC BY-SA. Présentation adaptée par Maths Explorer.}\par\medskip
\medskip

\begin{enumerate}
	\item $a = \ln (9) + \ln \left(\dfrac{\sqrt{3}}{3} \right)  + \ln \left(\dfrac19 \right)=\ln\left(3^2\right)+\ln\left(\sqrt{3}\right)-\ln(3) -\ln\left(3^2\right)=1\dfrac12\ln(3)-\ln(3)=-\dfrac12\ln(3)$
	
	\hfill \fbox{\textbf{Réponse d}}
	\item Donnons en premier le domaine de définition. Il faut que $x> 0$ et $x>10$. Donc 
	$\mathcal{D}_f = ]10~;~+\infty[$.
	
Dans $\mathcal{D}_f, \: \ln(x) + \ln (x - 10) = \ln(3) + \ln(7) \iff \ln\left(x(x - 10)\right)= \ln(3 \times 7) \iff x(x-10) = 21$  (par croissance de la fonction logarithme népérien) ou 

$x^2 - 10x - 21 = 0$. $\Delta=184 > 0$, donc ce trinôme admet deux solutions. Or $\sqrt{\Delta}=\sqrt{184} = \sqrt{4 \times 46} = 2\sqrt{46}$.

Les deux solution sont : $x_1= \dfrac{10 -2\sqrt{46}}{2}=5 -\sqrt{46}$ et $x_2=\dfrac{10+2\sqrt{46}}{2}=5+\sqrt{46}$.

Or $5^2 = 25 < 46$, donc $5 < \sqrt{46} \iff 5 - \sqrt{46}< 0$ et $5 - \sqrt{46} \notin \mathcal{D}_f$.

Par contre $\sqrt{46} > 5 \Rightarrow 5 + \sqrt{46} > 5 + 5 = 10$, donc $x_2 \in \mathcal{D}_f$.

La seule solution réelle est donc $5 + \sqrt{46}$.
\hfill \fbox{\textbf{Réponse c}}.
	 \item La fonction $f$ est définie sur l'intervalle $]0~;~+ \infty[$ par l'expression $f(x) = x^2(- 1 + \ln(x))$. Elle est continue et dérivable sur $]0~;~+ \infty[$.

	 $\forall x \in ]0~;~+ \infty[$, $f'(x)=2x\times (- 1 + \ln(x))+x^2\times \dfrac1x=-2x+2x\ln(x)+x=-x+2x\ln(x)=x(2\ln(x)-1)$.

	 $f'(x)=0 \iff x(2\ln(x)-1)=0 \iff 2\ln(x)-1 \iff \ln(x)=\dfrac{1}{2} \iff x=\e^{\frac12}=\sqrt{\e}$.
	 La tangente au point d'abscisse $a=\sqrt{\e}$ est horizontale. Son équation est donnée par : $y=f'(a)(x-a)+f(a)$.
	 Sachant que $f'\left(\sqrt{\e}\right)=0$ et que $f\left(\sqrt{\e}\right)= \sqrt{\e}^2\left(-1+\ln(\sqrt{\e})\right)=\e\left(-1+\dfrac12\right)= -\dfrac12\e$.
	 L'équation de cette tangente horizontale est donc : $y = -\dfrac12\e$.

\hfill \fbox{\textbf{Réponse d}}
	 \item Un sac contient 20 jetons jaunes et 30 jetons bleus. On tire successivement et avec remise 5 jetons du sac.

Il s'agit là d'un schéma de Bernoulli : la répétition de 5 expériences aléatoires n'ayant que deux issues, identiques et indépendantes entre elles. Si on note par $X$ la variable aléatoire donnant le nombre de jetons jaune tirés, alors $X$ suit donc une loi binomiale de paramètres $n = 5$ et $p = \dfrac{20}{50}=\dfrac25$ : $X \sim \mathcal{B}\left(5~;~\dfrac25\right)$

Ainsi $p(X=2)=\dbinom{5}{2}\times \left(\dfrac25\right)^2 \times \left(1-\dfrac25\right)^{5-2}=0,3456\approx 0,346$.

\hfill \fbox{\textbf{Réponse b}}
\item On reprend les conditions précédentes : un sac contient 20 jetons jaunes et 30 jetons bleus. On tire successivement et avec remise 5 jetons du sac. $X$ est la variable aléatoire donnant le nombre de jetons jaune tirés et $X$ suit donc une loi binomiale de paramètres $n=5$ et $p=\dfrac{20}{50}=\dfrac25$ : $X \sim \mathcal{B}\left(5~;~\dfrac25\right)$.

$p(X\geqslant 1)=1-p(X<1)=1-p(X=0)=1-\left(\dfrac35\right)^5\approx0,922$.

\hfill \fbox{\textbf{Réponse d}}
	\item On reprend les conditions précédentes : un sac contient 20 jetons jaunes et 30 jetons bleus. On tire successivement et avec remise 5 jetons du sac. $X$ est la variable aléatoire donnant le nombre de jetons jaune tirés et $X$ suit donc une loi binomiale de paramètres $n=5$ et $p=\dfrac{20}{50}=\dfrac25$ : $X \sim \mathcal{B}\left(5~;~\dfrac25\right)$.

L'espérance d'une loi binomiale est égale à : $E(X) = np$.

Ici, $E(x)=5\times \dfrac25=2$.

\hfill \fbox{\textbf{Réponse c}}
\end{enumerate}

\end{document}
